\section{Design space}
\label{sec:design}

The proposal has four separable decisions: (1) where the proofs are
\emph{checked} (custom kernel, proof assistant, SMT, or a tiered mix);
(2) where the proofs \emph{live} (comments, attributes, an embedded
script language, or an external proofs workspace keyed to the code);
(3) what the propositions are \emph{about} (a Rust memory model, an
ECMA-335 abstract machine, and the simulation between them); and (4)
how the pieces are \emph{compartmentalized} into independently
buildable, versionable, and abandonable packages. The decisions
interact, but they fail independently, so we treat them in turn.

\subsection{Option A: a bespoke dependently-typed checker}
\label{sec:design-custom}

The reading of the proposal closest to its letter: design a small
dependently-typed calculus --- $\Pi$/$\Sigma$ types, indexed inductive
families for the invariant taxonomy of \cref{sec:obligations},
definitional equality strong enough for layout arithmetic --- and
implement its kernel as a new workspace crate, with proof terms carried
beside the code and checked in CI.

The attraction is real: the calculus can be \emph{domain-fitted}.
Judgments like $\Gamma \vdash \pi : \rooted(o)$ or
$\Gamma \vdash \stw \Rightarrow \live(a)$ become primitive forms; proof
terms stay short because the type theory knows about arenas, layouts,
and phases. And a kernel for a bidirectional dependently-typed checker
is a known quantity --- a few thousand lines for the core.

The case against is decisive at this project's scale, and it is worth
being explicit because it answers the question as asked:

\begin{enumerate}[leftmargin=1.6em]
\item \textbf{The kernel joins the TCB with no independent audit.} A
  bespoke checker occupies the same trust position VeriFast occupied for
  fifteen years~\cite{foundationalverifast2026} --- ``checked by a tool
  whose soundness argument is on paper'' --- except VeriFast has had
  thousands of users shaking it. Lean and Rocq kernels are small,
  ancient, and independently re-implemented; nothing built here in a year
  approaches that.
\item \textbf{Automation is the actual cost center, not the kernel.}
  The corpus is dominated by obligations with heavy arithmetic and
  layout content (offset/alignment computations, bounds, descriptor
  interning). In Lean these fall to \code{omega}/\code{grind}/%
  \code{lean-smt}; in Rocq to \code{lia} and Iris's proof mode; in a
  bespoke system every one of these must be rebuilt or the proofs
  become unwritable. The history of ACSL~\cite{acsl} says the SMT-for-%
  the-bulk / assistant-for-the-residue split is what makes annotation
  burden bearable; a bespoke kernel starts with neither half.
\item \textbf{The propositions still need a model.} Even a perfect
  kernel proves theorems only relative to axioms about Rust and the VES.
  Those axiomatizations are the hard, novel work
  (\cref{sec:design-model}); they are needed under every option, and
  under Option A they get the least scrutiny.
\item \textbf{It forfeits the ecosystem the workflow depends on.} The
  author's stated workflow is LLM-agent-driven. Lean 4 has the largest
  machine-proof training corpus and tooling surface of any assistant;
  a private calculus has zero. The single biggest practical lever ---
  agents drafting and repairing proofs under a checker's supervision ---
  works dramatically better against a mainstream target.
\end{enumerate}

\begin{riskbox}[Assessment of Option A]
Building the calculus is feasible; trusting it is not justifiable,
and its automation alone would exceed the budget of a reasonable
first campaign. A domain-fitted \emph{surface language} is
nonetheless the right idea --- \cref{sec:dsl} designs one, including a
tactic layer --- but as untrusted syntax elaborated into an existing
kernel's terms, never as a kernel. The tactic point deserves
emphasis, because it is where the bespoke-kernel temptation recurs:
domain tactics (``open the STW invariant,'' ``frame this token'') are
implemented as Lean metaprograms whose \emph{output} the kernel
checks, so the surface can grow arbitrarily specialized while the TCB
does not grow at all.
\end{riskbox}

\subsection{Option B: an embedded script DSL elaborated into Lean 4}
\label{sec:design-embed}

Here the DSL is a block-shaped, tactic-oriented proof script embedded
in the Rust source as a macro that \emph{wraps} each \code{unsafe}
site --- statement-shaped lines in the lineage of \code{dotnetdll}'s
\code{asm!} macro~\cite{dotnetdll} (\cref{sec:dsl-lineage}) --- naming
a goal (declared or inferred), acquiring premise witnesses as
zero-sized Rust \emph{ghost tokens} the compiler itself checks, and
sketching the discharge with vocabulary tactics. The scripts elaborate
deterministically into ordinary Lean~4 statements in a proofs
workspace, where the proofs are ordinary Lean terms and the checker is
Lean's kernel. A small external tool bridges the two worlds by
comparing content-hashed manifests --- never by running Lean inside the
VM's build (\cref{sec:design-compart}).

This is the architecture this study recommends, and the backend choice
follows from \cref{sec:prior-backends}:

\begin{itemize}[leftmargin=1.4em]
\item \textbf{Lean 4 as the home} for the ECMA-335 abstract machine
  (executable, JSCert-style, differentially tested against the existing
  fixture oracle), for the invariant-family library and its tactic
  layer, and for the simulation-relation theorems. Lean's
  metaprogramming carries both ends of the DSL: the vocabulary is
  compiled from one declarative source into Rust token types and Lean
  declarations, and the surface tactics are Lean tactic macros.
\item \textbf{Rust-substrate obligations stated against a trusted
  axiomatization} of the MiniRust/Tree-Borrows UB rules --- a Lean
  library of axioms (or definitions validated by testing) rather than a
  from-scratch re-mechanization of what already exists in Rocq. The
  alternative --- doing the whole project in Rocq to sit on
  Iris/RustBelt/Tree Borrows --- buys foundational depth on the Rust side
  at the cost of the model half (no executable-spec tradition as strong,
  weaker agent tooling, and the VES model is the bigger artifact).
  \Cref{sec:risks} treats the residual gap honestly: theorems about the
  Rust side are then conditional on the axiomatization's fidelity.
\item \textbf{An escape hatch upward}: for a handful of critical modules
  (the \gcarena{} boundary, the STW protocol), RefinedRust in Rocq
  remains available as a second campaign; nothing in the Lean-first
  design blocks it, because the script layer is checker-agnostic by
  construction (the goal names the obligation; which backend discharges
  it is a per-obligation choice recorded in the proof manifest).
\end{itemize}

\subsection{Option C: adopt an existing verifier wholesale}
\label{sec:design-adopt}

For completeness: pointing RefinedRust, VeriFast, or Verus at
\crate{dotnet-vm} as-is.

\begin{itemize}[leftmargin=1.4em]
\item \emph{RefinedRust} is the right theorem but the wrong scale today:
  partial language coverage meets a 26-crate workspace built on
  \gcarena{}'s branded lifetimes, GATs, and heavy trait towers; and the
  ECMA-335 layer would still have to be built inside its Rocq setting.
\item \emph{VeriFast} comes closest in form factor (proofs in comments,
  modules with \code{unsafe}, Tree Borrows obligations) and is the most
  credible off-the-shelf pilot: worth a two-week spike on
  \crate{dotnet-utils}' atomic-access module. Its separation logic knows
  nothing of the VES, so spec-level invariants become user-defined
  predicates --- workable, unassisted.
\item \emph{Verus} would mean rewriting the heap and stack cores against
  its permission tokens --- a different project (arguably a different,
  interesting project: a Verus-native VES), not a verification layer on
  this one. Its ghost-permission idiom is, however, the direct
  inspiration for the ghost-token layer of \cref{sec:dsl-tokens},
  reused here at the annotation layer rather than as a rewrite.
\end{itemize}

Option C fails as the whole answer for the same reason it succeeds as a
component: these tools verify Rust against Rust's model. The author's
actual ambition --- ``take our ECMA-335 target, combine it with our Rust
substrate, and show compliance'' --- requires the two-level structure
that none of them ships.

\subsection{The two-level model architecture}
\label{sec:design-model}

Whatever discharges the proofs, the propositions need the following
mechanized stack; this is the load-bearing engineering of the whole
proposal.

\begin{enumerate}[leftmargin=1.6em]
\item \textbf{$\mathcal{M}_{\mathrm{VES}}$ --- an executable ECMA-335
  abstract machine in Lean 4}, covering the invariant-relevant core: the
  typed memory store with alignment/atomicity rules
  (\ecma{I.12.6.2}, \ecma{I.12.6.6}--\ecma{I.12.6.8}), object/array
  layout including the reference-overlap prohibition (\ecma{II.10.7},
  \ecma{II.22.8}), the managed/unmanaged pointer taxonomy and
  GC-reporting rules (\ecma{II.14.4}, \ecma{III.1.1.5}), evaluation-stack
  typing for the executed CIL fragment (after
  Gordon--Syme~\cite{gordonsyme2001}), reachability/finalization
  (\ecma{I.8.9.6.7}), and --- as an implementation-level module, since the
  standard only licenses it --- the STW/arena protocol as a labelled
  transition system. Executability is non-negotiable: it is what lets
  the existing differential fixture harness test the \emph{model}
  against real .NET, JSCert-style~\cite{jscert2014}, before any proof
  depends on it.
\item \textbf{$\mathcal{M}_{\mathrm{Rust}}$ --- a shallow Lean fragment
  of the Rust value/memory model} sufficient to state UB side conditions:
  typed byte ranges, provenance-carrying pointers, validity/initialization
  predicates, data-race conditions, and the Tree Borrows obligations for
  the retag patterns the codebase actually uses (a deliberately small,
  test-validated axiomatization per \cref{sec:design-embed}).
\item \textbf{$\mathcal{R}$ --- the representation relation} connecting
  the two: layout descriptors computed by \code{LayoutManager} correspond
  to $\mathcal{M}_{\mathrm{VES}}$ layouts; \code{ObjectRef}/%
  \code{ObjectPtr}/\code{ManagedPtr} values represent spec-level
  references; arena/coordinator states refine the protocol LTS. Each
  invariant family of \cref{sec:obligations} becomes a predicate stated
  once over $\mathcal{R}$ and instantiated by many \safetyc{} sites.
\end{enumerate}

The payoff structure matters: $\mathcal{M}_{\mathrm{VES}}$ is useful
\emph{before any unsafe block is verified} --- as a differential oracle
and as the first mechanized ECMA-335 model, a publishable artifact in
its own right. The proposal should be staged so that this asset lands
first (\cref{sec:roadmap}).

\subsection{Compartmentalization: packages, boundaries, artifacts}
\label{sec:design-compart}

The proposal's suggestion to spin the verifier into its own workspace
is realized as two sibling repositories: \code{ves-proof}, whose Rust
packages live under \code{crates/}, and \code{ves-proof-lean}, whose
independent Lake packages live under \code{packages/}.  This concrete
split generalizes to the design rule on which the project's viability
rests: \emph{every piece is a separately versioned package; the only
things crossing package boundaries are code dependencies pointing in
one direction or serialized artifacts with versioned schemas; and
\dotnetrs{}'s ordinary build never invokes Lean.}
\Cref{tab:packages} lays out the decomposition;
\cref{fig:dependencies} the dependency direction.

\begin{table}
  \centering
  \small
  \begin{tabular}{@{}lllp{13.2em}@{}}
    \toprule
    Package & Kind & CI owner & Role and independence \\
    \midrule
    \code{ves-syntax}     & Rust lib        & \code{ves-proof} & script grammar, AST, canonical serializer, hashing; no proc-macro, no Lean; the \emph{single} parser linked by both macro and checker \\
    \code{ves-vocabulary} & data + codegen  & \code{ves-proof} & declarative goal/premise/tactic source; generates \code{ves-tokens} types \emph{and} \code{VesCore} declarations; semver'd; the compatibility anchor \\
    \code{ves-tokens}     & Rust lib (tiny) & \code{ves-proof} & ZST ghost tokens, cfg-mirrored constructors, \code{TrustGrant}; \code{no\_std}, zero-cost; the only runtime dependency \dotnetrs{} gains \\
    \code{ves-macros}     & proc-macro      & \code{ves-proof} & \code{proof!}, \code{proof\_impl!}, \code{contract!}; parsing via \code{ves-syntax}; expansion = tokens + verbatim subject \\
    \code{ves-check}      & CLI bin         & \code{ves-proof} & extraction, structural rules, manifest diff, ratchet budgets; needs no Lean toolchain \\
    \midrule
    \code{VesCore}        & Lean (Lake)     & \code{ves-proof-lean} & vocabulary mirror, goal elaborator, statement hashing, \code{ves\_sketch} tactic frame \\
    \code{VesModel}       & Lean (Lake)     & \code{ves-proof-lean} & $\mathcal{M}_{\mathrm{VES}}$, executable; differential test harness against .NET fixtures \\
    \code{RustAssumptions}& Lean (Lake)     & \code{ves-proof-lean} & $\mathcal{M}_{\mathrm{Rust}}$ axioms + Miri-falsifier index \\
    \code{VesProtocol}    & Lean (Lake)     & \code{ves-proof-lean} & STW/arena LTS, protocol invariants and theorems \\
    \code{DotnetRsProofs} & Lean (Lake)     & \code{ves-proof-lean} & $\mathcal{R}$ + site lemmas; exports the proof manifest \\
    \midrule
    \dotnetrs{} integration & workspace changes & \dotnetrs{} & scripts at unsafe sites; dev-deps on \code{ves-macros}/\code{ves-tokens}; \code{xtask ves-check} CI job; \code{TRUSTED.toml} \\
    \bottomrule
  \end{tabular}
  \caption{The package decomposition. The five Rust crates exist under
  \code{ves-proof/crates/}; the five Lean packages exist under
  \code{ves-proof-lean/packages/}.  Each is publishable and testable in
  isolation.}
  \label{tab:packages}
\end{table}

\begin{figure}
\centering
\begin{lstlisting}[language={},frame=none,basicstyle=\ttfamily\footnotesize]
                 (code dependencies: arrow = "depends on")

  dotnet-rs ----> ves-macros ----> ves-syntax <---- ves-check
      |               |                                  |
      +-------------> ves-tokens <--- [generated] ---+   | (CI tool,
                                                     |   |  not a dep)
              ves-vocabulary --[codegen]--> VesCore  |
                                              ^      |
        DotnetRsProofs -> VesProtocol ---------+------+
             |       \--> RustAssumptions -> VesCore
             v
          VesModel

                 (artifact flow: data, never code)

  dotnet-rs  ==obligations.v1.json==>  DotnetRsProofs (stub gen)
  DotnetRsProofs ==proofs.v1.json==>   dotnet-rs CI (ves-check diff)
  TRUSTED.toml   ==TrustGrant codegen==> dotnet-rs build
\end{lstlisting}
\caption{Dependency direction and artifact flow. No Lean package
depends on \dotnetrs{} source; \dotnetrs{} depends on no Lean package;
the two meet only in serialized manifests.}
\label{fig:dependencies}
\end{figure}

\paragraph{The two artifact schemas.} Boundaries are serialization
contracts, versioned like any other API:

\begin{itemize}[leftmargin=1.4em]
\item \emph{Obligation manifest} (\code{obligations.v1.json}; producer
  \code{ves-check extract}, consumers \code{VesCore}'s stub generator
  and the structural rules): per \code{(site, cfg-set)} --- goal term,
  premise list with trust classes, tactic sketch, source span,
  statement hash, subject hash, vocabulary and grammar versions.
  Canonical JSON (sorted keys, normalized paths, no timestamps) so the
  same tree always serializes byte-identically.
\item \emph{Proof manifest} (\code{proofs.v1.json}; producer
  \code{DotnetRsProofs}' Lake export, consumer \code{ves-check} in
  \dotnetrs{} CI): obligation hash $\to$ theorem name,
  \code{sorry}-free flag, transitive axiom set, Lean toolchain,
  vocabulary version. Published as a tagged release artifact;
  \dotnetrs{} pins a version exactly as it pins a crate.
\end{itemize}

\paragraph{Compatibility rules.} The vocabulary version is the
compatibility anchor: statement hashes embed it, both manifests carry
it, and \code{ves-check} refuses to compare manifests across a
vocabulary major version. Additive vocabulary changes (new goal
constructors, new tactics) are minor and never invalidate existing
hashes; changing a constructor's elaboration is major by definition
because it moves hashes. Proof manifests may lag the code: an
obligation with no matching proof entry is \emph{unproved} (ratcheted,
warning-level by default; blocking for sites the policy designates
critical),
while structural violations --- unwrapped \code{unsafe}, unregistered
trust grants, budget breaches --- always block. This split is what
lets \dotnetrs{} refactor at full speed against a week-old proof
manifest without ever weakening the structural floor.

\paragraph{CI ownership.} \dotnetrs{}'s CI gains one fast job: build
(which already type-checks every ghost token in every matrix leg) plus
\code{ves-check} against the pinned proof manifest --- seconds, no
Lean. The \code{ves-proof} repository owns Rust tooling and its tests;
\code{ves-proof-lean} owns Lake builds, the \code{VesModel}
differential suite against real .NET, and the Miri falsifier index for
\code{RustAssumptions}.  A scheduled verifier job may coordinate both
repositories, pull \dotnetrs{} HEAD, re-extract obligations, and open a
drift report when hashes have moved --- the cross-repository analogue
of \code{check\_doc\_drift.sh}.

\paragraph{Independence and abandonability.} Each layer is useful if
everything above it stops: vocabulary + macros + tokens alone give
rustc-checked premise discipline with zero Lean (Phase 1);
\code{ves-check} with an empty proof manifest is a
tag-std-class~\cite{tagstd2025} structural layer; \code{VesModel} is a
publishable mechanized ECMA-335 core with no dependence on the DSL at
all; \code{VesProtocol} and \code{RustAssumptions} are independently
meaningful theories; only \code{DotnetRsProofs} is coupled to the VM,
and abandoning it merely freezes every site at \code{trusted} ---
strictly better than today's prose, at no ongoing cost. Conversely,
nothing in \dotnetrs{} breaks if either verifier repository goes quiet:
the pinned Rust crates still expand and type-check, and the pinned
manifest still describes exactly what was once proved.
